% =========================================================================
% Seminar / Conference Talk Template
%
% Scaffold for /create-talk. Rules live in style.md + structure.md
% (loaded with this skill); fidelity rules in single-source-of-truth.md;
% past-talk pitfalls in quality_reports/2026-05-04_seminar_talk_review.md.
%
% Format budgets (slides at 60-90 sec/slide pacing):
%   30-min seminar       25-32 main + 10-15 backup
%   45-60-min job market 40-50 main + 15-20 backup  (add Preview after Block 3)
%   15-min short         10-15 main, few backup
%   5-min lightning       3-5 main, no backup
%
% slides.sty macros (most-used):
%   \bi/\ei \be/\ee     itemize / enumerate
%   \alert{} \asher{}   raspberry emphasis / gray de-emphasis
%   \xsmallcitecolor{\citet{...}}    grayed inline citation
%   \highlight{...}     equation-term wrap (raspberry); also \highlightteal/orange
%   \tabletitle{Outcome:}{...}       title OUTSIDE the tabular
%   \bottomleft{\beamergotobutton{Back}}    appendix pills / back-buttons
%   \imageframe{path}   full-bleed image frame
%
% Colors defined in slides.sty:
%   raspberry teal navy trueblue digitalblue lightblue
%   purple cranberry orange ruby alice daisy coral kelly asher (gray)
%   accent / accent2 (\providecolor — do not redefine here)
%
% Replace every [TODO: ...] before delivery.
% =========================================================================

\documentclass[aspectratio=169,t,11pt,table]{beamer}
\usepackage{slides,math}

\graphicspath{{figures/}{../}}

\title{[TODO: paper title]}
\author{[TODO: author names]}
\institute{[TODO: affiliation] \\ [TODO: conference / seminar / lecture name]}
\date{[TODO: Month YYYY]}

\begin{document}

% BLOCK 1 — TITLE
\begin{frame}[noframenumbering,plain]
  \maketitle
\end{frame}

% BLOCK 2 — HOOK / FRAMING (2-4 slides)
% Field-level puzzle. No paper-specific machinery, no findings, no
% citations on the first 1-2 slides.
\begin{frame}{[TODO: declarative phenomenon — e.g., "Decisions increasingly rely on algorithmic prediction"]}
  % TODO: open the field-level question. Avoid jargon.
  \pause
  \bi
    \item [TODO: 1st piece of context — concrete, layperson-accessible]
    \item [TODO: 2nd piece of context]
  \ei
\end{frame}

\begin{frame}{[TODO: framing slide 2 — sharpen the puzzle]}
  % TODO: introduce the tension the paper resolves
\end{frame}

% BLOCK 3 — PAPER INTRODUCTION (1 slide, 3 builds)
% Build 1: questions. Build 2: setting/technology/design labels.
% Build 3: \alert{} answers. Co-author line in second \frame{}{} arg.
\begin{frame}{A paper: [TODO: short title]}{[TODO: Joint work with X (Affiliation) and Y (Affiliation)]}
  \be
    \item<1-> [TODO: question 1 — e.g., "Does X cause Y?"]
      \only<3>{ \\ \alert{[TODO: 1-line answer with magnitude]} }
    \item<1-> [TODO: question 2]
      \only<3>{ \\ \alert{[TODO: 1-line answer with magnitude]} }
    \item<1-> [TODO: question 3]
      \only<3>{ \\ \alert{[TODO: 1-line answer with magnitude]} }
  \ee
  \onslide<2->{
    \medskip
    \textcolor{teal}{Setting:} [TODO: 1-line setting] \\
    \textcolor{teal}{Technology / Treatment:} [TODO: 1-line] \\
    \textcolor{teal}{Empirical design:} [TODO: 1-line]
  }
\end{frame}

% BLOCK 4 — LITERATURE POSITIONING (1 slide)
% "Optimism and apprehension" pattern. Two strands as builds. Final
% build: "This paper:" block. Placement is flexible (before or after
% Block 3 depending on whether the literature debate IS the puzzle).
\begin{frame}{[TODO: Topic]: optimism and apprehension}
  \textcolor{trueblue}{Optimism:}
  \bi
    \item<1-> [TODO: claim 1] \xsmallcitecolor{\citet{[TODO: ref1, ref2]}}
    \item<1-> [TODO: claim 2] \xsmallcitecolor{\citet{[TODO: ref3]}}
  \ei
  \pause
  \medskip
  \textcolor{raspberry}{Apprehension:}
  \bi
    \item<2-> [TODO: counter-claim 1] \xsmallcitecolor{\citet{[TODO: ref4, ref5]}}
    \item<2-> [TODO: counter-claim 2] \xsmallcitecolor{\citet{[TODO: ref6]}}
  \ei
  \pause
  \bigskip
  \alert{This paper:} [TODO: 1-2 sentences naming the gap and the contribution]
\end{frame}

% BLOCK 5 — AGENDA (1 slide, no builds)
% Re-shown at each section transition; completed sections in \asher{}.
\begin{frame}{Agenda}
  \be
    \item Setting \& data
    \item Empirical strategy
    \item [TODO: Section 3 title — first finding block]
    \item [TODO: Section 4 title — framework + second finding block]
    \item [TODO: Section 5 title — heterogeneity / robustness]
  \ee
\end{frame}

% BLOCK 6 — SETTING (3-5 slides)
% Pattern: why-this-setting → concrete-example → why-X-is-an-ML-problem.
% Drop or rename the third slide for non-ML papers. Add a 4th for the
% treatment object (AI tool, policy, etc.) when relevant.
\begin{frame}{Why [TODO: setting]?}
  \bi
    \item<1-> \alert{[TODO: reason 1 keyword]} — [TODO: 1-line elaboration]
    \item<2-> \alert{[TODO: reason 2 keyword]} — [TODO: 1-line elaboration]
    \item<3-> \alert{[TODO: reason 3 keyword]} — [TODO: 1-line elaboration]
  \ei
\end{frame}

\begin{frame}{[TODO: concrete example title]}
  % Use \begin{quotation}...\end{quotation} for an interview quote.
  % Citation underneath in raspberry or trueblue.
  \begin{center}
    \includegraphics[width=0.7\textwidth]{[TODO: screenshot.png]}
  \end{center}
\end{frame}

\begin{frame}{Why is [TODO: setting] an [TODO: ML / prediction] problem?}
  \be
    \item<1-> \alert{[TODO: prediction structure 1]} — [TODO: 1-line]
    \item<2-> \alert{[TODO: prediction structure 2]} — [TODO: 1-line]
    \item<3-> \alert{[TODO: prediction structure 3]} — [TODO: 1-line]
  \ee
\end{frame}

% BLOCK 7 — DATA (1-2 slides)
% Three labeled blocks (Sample / Unit info / Outcomes); pills to backup.
\begin{frame}{Our data}
  \textcolor{raspberry}{Sample:} [TODO: 1-line — N units, time range] \\[2pt]
  \textcolor{teal}{[TODO: Unit]-level information:}
  \bi
    \item [TODO: variable type 1]
    \item [TODO: variable type 2]
  \ei
  \medskip
  \textcolor{digitalblue}{Outcomes:}
  \bi
    \item [TODO: outcome 1 (with unit)]
    \item [TODO: outcome 2 (with unit)]
  \ei
  \bottomleft{
    \hyperlink{summary-stats}{\beamergotobutton{Summary stats}}
    \hyperlink{balance-table}{\beamergotobutton{Balance}}
  }
\end{frame}

% BLOCK 8 — EMPIRICAL STRATEGY (1-3 slides)
% Visual-first, then equation walkthrough with 4-5 term-by-term builds
% (dep var → FE → treatment → controls → plain + assumptions).
% IV: repeat the equation slide for first stage + reduced form.
% Triple-diff / structural: one slide per equation tier.
\begin{frame}{Empirical strategy: [TODO: 1-line description of variation]}
  \begin{columns}
    \begin{column}{0.6\textwidth}
      \includegraphics[width=\linewidth]{[TODO: rollout_figure.png]}
    \end{column}
    \begin{column}{0.4\textwidth}
      \bi
        \item [TODO: constraint / source of variation]
        \item [TODO: identifying assumption in plain English]
      \ei
    \end{column}
  \end{columns}
\end{frame}

\begin{frame}{Empirical strategy: equation}
  \only<1>{
    $$ \textcolor{raspberry}{\underbrace{\bm{y_{it}}}_{\text{[TODO: outcome]}}}
       = \alpha_i + \delta_t + \beta D_{it} + \gamma X_{it} + \varepsilon_{it} $$
  }
  \only<2>{
    $$ y_{it}
       = \textcolor{raspberry}{\underbrace{\bm{\alpha_i + \delta_t}}_{\text{[TODO: FE structure]}}}
       + \beta D_{it} + \gamma X_{it} + \varepsilon_{it} $$
  }
  \only<3>{
    $$ y_{it} = \alpha_i + \delta_t
       + \textcolor{raspberry}{\underbrace{\bm{\beta D_{it}}}_{\text{[TODO: treatment]}}}
       + \gamma X_{it} + \varepsilon_{it} $$
  }
  \only<4>{
    $$ y_{it} = \alpha_i + \delta_t + \beta D_{it}
       + \textcolor{raspberry}{\underbrace{\bm{\gamma X_{it}}}_{\text{[TODO: controls]}}}
       + \varepsilon_{it} $$
  }
  \only<5>{
    $$ y_{it} = \alpha_i + \delta_t + \beta D_{it} + \gamma X_{it} + \varepsilon_{it} $$
    \medskip
    \bi
      \item \alert{Identifying assumption:} [TODO]
      \item \alert{Cluster standard errors at:} [TODO]
    \ei
  }
\end{frame}

% AGENDA RE-SHOW before Findings
\begin{frame}{Agenda}
  \be
    \item \asher{Setting \& data}
    \item \asher{Empirical strategy}
    \item \textbf{[TODO: Section 3 title]}
    \item \asher{[TODO: Section 4 title]}
    \item \asher{[TODO: Section 5 title]}
  \ee
\end{frame}

% BLOCK 9 — FINDINGS, FIRST HALF (4-8 slides)
% Declarative title (state the finding, not "Main results"). Figure body
% with progressive builds; optional magnitude sidebar; pills to backup.
% Standard follow-up: "Similar results across estimators".
\begin{frame}{[TODO: headline finding — declarative, with magnitude]}
  \begin{columns}
    \begin{column}{0.7\textwidth}
      \includegraphics[width=\linewidth]{[TODO: main_finding.png]}
    \end{column}
    \begin{column}{0.3\textwidth}
      \alert{Magnitude:} [TODO: % effect or units] \\[6pt]
      \alert{Sample:} [TODO: N units] \\[6pt]
      [TODO: 1-line takeaway]
    \end{column}
  \end{columns}
  \bottomleft{
    \hyperlink{dd-table}{\beamergotobutton{Table}}
    \hyperlink{altspecs}{\beamergotobutton{Alt. specifications}}
  }
\end{frame}

\begin{frame}{Similar results across estimators}
  \begin{center}
    \includegraphics[width=0.8\textwidth]{[TODO: estimator_comparison.png]}
  \end{center}
\end{frame}

% Repeat finding slides as needed for first half (typically 2-4 more)

% AGENDA RE-SHOW before Framework
\begin{frame}{Agenda}
  \be
    \item \asher{Setting \& data}
    \item \asher{Empirical strategy}
    \item \asher{[TODO: Section 3 title]}
    \item \textbf{[TODO: Section 4 title]}
    \item \asher{[TODO: Section 5 title]}
  \ee
\end{frame}

% BLOCK 10 — CONCEPTUAL FRAMEWORK (2-3 slides, MID-DECK)
% Place AFTER the first finding. Two-color comparison (trueblue vs
% raspberry), 3 builds: defining bullets → parallel equations → diff.
\begin{frame}{[TODO: Framework title — e.g., "Humans and algorithms are distinct prediction technologies"]}
  \begin{columns}
    \begin{column}{0.5\textwidth}
      \centering
      \textcolor{trueblue}{[TODO: Entity A]} \\
      {\color{trueblue}\rule{\linewidth}{2pt}}
      \bi
        \item<1-> \textcolor{trueblue}{[TODO: defining property]}
        \item<2-> $\hat{y}^A_i = [TODO]$
      \ei
    \end{column}
    \begin{column}{0.5\textwidth}
      \centering
      \textcolor{raspberry}{[TODO: Entity B]} \\
      {\color{raspberry}\rule{\linewidth}{2pt}}
      \bi
        \item<1-> \textcolor{raspberry}{[TODO: defining property]}
        \item<2-> $\hat{y}^B_i = [TODO]$
      \ei
    \end{column}
  \end{columns}
  \onslide<3->{
    \bigskip
    \centering
    $\hat{y}^A_i - \hat{y}^B_i = [TODO: difference equation]$
    \medskip
    \alert{Prediction:} [TODO: testable implication]
  }
\end{frame}

% BLOCK 11 — FINDINGS, SECOND HALF (4-6 slides)
% Same template as Block 9. Tests framework predictions, not headline.
\begin{frame}{[TODO: sub-finding 1 — declarative]}
  \begin{center}
    \includegraphics[width=0.8\textwidth]{[TODO: subfinding1.png]}
  \end{center}
\end{frame}

\begin{frame}{[TODO: heterogeneity finding — declarative]}
  \begin{center}
    \includegraphics[width=0.8\textwidth]{[TODO: heterogeneity.png]}
  \end{center}
\end{frame}

% BLOCK 12 — ROBUSTNESS (2-4 slides)
% 3+ challenges: checklist re-show with progressive graying.
% 1-2 challenges: direct deep-dive, skip the re-shows.
\begin{frame}{Assessing the empirical strategy}
  \bi
    \item<1-> \textbf{Challenge 1:} [TODO: e.g., "Pre-trends"]
    \item<2-> \textbf{Challenge 2:} [TODO: e.g., "Selection into treatment"]
    \item<3-> \textbf{Challenge 3:} [TODO: e.g., "Spillovers"]
  \ei
\end{frame}

\begin{frame}{[TODO: deep-dive on Challenge 1]}
  \begin{center}
    \includegraphics[width=0.7\textwidth]{[TODO: robustness1.png]}
  \end{center}
\end{frame}

\begin{frame}{Assessing the empirical strategy}
  \bi
    \item \asher{Challenge 1: [TODO]} \checkmark
    \item \textbf{Challenge 2:} [TODO]
    \item Challenge 3: [TODO]
  \ei
\end{frame}

% CONCLUSION (1-2 slides)
% "When [trigger]:" → enumerated findings → \alert{} so-what → email.
\begin{frame}{Conclusion}
  \textbf{When [TODO: trigger condition mirroring research question]:}
  \be
    \item [TODO: Finding 1 with magnitude]
    \item [TODO: Finding 2 with magnitude]
    \item [TODO: Finding 3 with magnitude]
  \ee
  \medskip
  \alert{[TODO: One-line so-what claim]:}
  \be
    \item [TODO: Implication 1]
    \item [TODO: Implication 2]
  \ee
  \medskip
  \textcolor{raspberry}{[TODO: email@institution.edu]}
\end{frame}

% APPENDIX + BACKUP (8-15 slides)
% Each backup frame: \label{name} matching the originator's \hyperlink,
% single content (no builds), tiny italic methodology note, back-button.
\appendix
\begin{frame}[allowframebreaks,noframenumbering]{Appendix}\end{frame}

\begin{frame}{Summary statistics}\label{summary-stats}
  \begin{center}
    \scalebox{.75}{\makebox[\linewidth]{\input{tables/[TODO: summary_stats.tex]}}}
  \end{center}
  \parbox{0.8\textwidth}{\justifying \tiny \textit{Note: [TODO: methodology note].}}
  \bottomleft{\hyperlink{data}{\beamergotobutton{Back}}}
\end{frame}

\begin{frame}{Balance table}\label{balance-table}
  \begin{center}
    \scalebox{.75}{\makebox[\linewidth]{\input{tables/[TODO: balance.tex]}}}
  \end{center}
  \parbox{0.8\textwidth}{\justifying \tiny \textit{Note: [TODO: methodology note].}}
  \bottomleft{\hyperlink{data}{\beamergotobutton{Back}}}
\end{frame}

\begin{frame}{[TODO: backup table title]}\label{dd-table}
  \tabletitle{Outcome:}{[TODO: $Y_{it}$ description]}
  \begin{center}
    \scalebox{.7}{\makebox[\linewidth]{\input{tables/[TODO: main_table.tex]}}}
  \end{center}
  \bottomleft{\hyperlink{[TODO: main finding label]}{\beamergotobutton{Back}}}
\end{frame}

\begin{frame}{Alternative specifications}\label{altspecs}
  \begin{center}
    \scalebox{.7}{\makebox[\linewidth]{\input{tables/[TODO: altspecs.tex]}}}
  \end{center}
  \bottomleft{\hyperlink{[TODO]}{\beamergotobutton{Back}}}
\end{frame}

\end{document}
